\section{The finite unitary character associated with the coefficients}
\label{sec:fd-unitary}
This section gives the exact change of variables that turns coefficient positivity into positivity on a fixed unitary group.
For \(a\ge0\) put
\[
 p(a)=F_a'(0),\qquad E_a(w)=F_a(w/p(a)),\qquad
 \Psi_a(z)=E_a(z-1)=\sum_{k\ge0}a_k(a)z^k.
\]
Thus \(E_a(0)=E_a'(0)=1\) and \(\Psi_a(1)=1\).
For a unitary \(d\) by \(d\) matrix \(U\) with eigenvalues
\(z_1,\ldots,z_d\), define
\[
 \chi_{a,d}(U)=\prod_{\nu=1}^d\Psi_a(z_\nu).
\]
A normalized character here means a continuous central function of
positive type with value one at the identity: for every finite list
\(U_i\), its matrix \([\chi(U_i^{-1}U_j)]_{i,j}\) is positive
semidefinite.

\begin{proposition}[Positivity at every fixed unitary rank]
\label{prop:fd-unitary}
For each \(d\ge1\) there is \(\eta_d>0\) such that \(\chi_{a,d}\)
is a normalized character for \(0\le a\le\eta_d\). Its expansion is
\[
 \chi_{a,d}(U)=\sum_{\ell(\lambda)\le d}
   c_{\lambda,d}(a)s_\lambda(U),\qquad
 c_{\lambda,d}(a)=\det[a_{\lambda_i-i+j}(a)]_{i,j=1}^d>0.
\]
Here \(\lambda\) is a partition padded with zeros, \(s_\lambda\)
is its irreducible polynomial character, and \(a_k=0\) for \(k<0\).
The expansion converges absolutely and uniformly, and
\[
 \sum_{\ell(\lambda)\le d}c_{\lambda,d}(a)s_\lambda(1^d)=1.
\]
For each \(a>0\), some finite rank fails positive type. If \(d_*(-a)\)
is the first such rank, then \(d_*(-a)\to\infty\) as \(a\downarrow0\).
Near zero it is exactly the least negative Toeplitz-minor order
of the translated coefficients \((a_k(a))_{k\ge0}\).
\end{proposition}

\begin{proof}
The genus-zero product takes the form
\[
 E_a(w)=\prod_j(1+\beta_j(a)w),\qquad
 \beta_j(a)=\frac1{p(a)z_j(-a)^2},\qquad
 \sum_j\beta_j(a)=1.
\]
There is one parameter for each right-half-plane heat zero, with
multiplicity. At \(a=0\) these parameters are positive and there
are infinitely many. Their sum is one, so their maximum is
strictly less than one. Choose a fixed \(b<1\) larger than that
maximum. The function \(E_0\) has no zero on \(|w|\le1/b\);
compact convergence gives the same property for small \(a\).
Thus every \(|\beta_j(a)|<b\).

Translation now gives an exact parameter map:
\[
 \Psi_a(z)=E_a(-1)\prod_j(1+\alpha_j(a)z),\qquad
 \alpha_j(a)=\frac{\beta_j(a)}{1-\beta_j(a)}.
\]
The real parameters \(\beta_j\) are positive near zero, their
factors \(1-\beta_j\) are positive, and conjugate parameters
give the positive product \(|1-\beta_j|^2\). Therefore
\(E_a(-1)>0\). The \(\alpha_j\) are summable in absolute value.
For \(\beta=re^{i\theta}\), \(r\le b\), differentiation of the
argument along the circular arc gives
\[
 \left|\arg\frac{\beta}{1-\beta}\right|
 \le\frac{|\theta|}{1-b}.
\]
Indeed the derivative of \(\log(\beta/(1-\beta))\) with respect
to \(\theta\) is \(i/(1-re^{i\theta})\), of modulus at most
\((1-b)^{-1}\). The angle at \(\theta=0\) is zero.
The uniform heat-angle estimate makes these translated angles
tend uniformly to zero. For each fixed \(d\), the pairwise
Schur criterion~\eqref{eq:fd-sector} and Cauchy--Binet now give
\(PF_d\) for \((a_k(a))\) on a small closed interval. The
positive path term proves strict positivity at every allowed
index, since there are still infinitely many factors.

We give the precise character calculation. For any entire
\(\Psi(z)=\sum a_kz^k\), multiply \(\prod_i\Psi(z_i)\) by the
Weyl alternant \(\det[z_i^{d-j}]\). Cauchy--Binet expands the
resulting determinant in alternants with decreasing exponents
\(\lambda_j+d-j\). Division by the original alternant gives
\[
 \prod_i\Psi(z_i)=
 \sum_{\ell(\lambda)\le d}
       \det[a_{\lambda_i-i+j}]_{i,j=1}^d s_\lambda(z_1,\ldots,z_d).
\]
The Weyl interpretation is
\cite[Theorem~5.22.1, p.~129]{Etingof2011}; the compact-group
positive-type convention is \cite[Section~4.1, p.~56]{Faraut2007}.
This identity is initially a degree-wise identity of power series.
To justify it uniformly on the unit torus, choose \(R>1\) and
let \(M_R=\max_{|z|=R}|\Psi(z)|\). Cauchy's estimate gives
\[
 \left|\det[a_{\lambda_i-i+j}]\right|
 \le d!M_R^d R^{-|\lambda|}.
\]
The sum of the indices in each determinant term is \(|\lambda|\);
terms with negative indices vanish. The Weyl dimension formula
bounds \(s_\lambda(1^d)\) by a polynomial in \(|\lambda|\), and
there are at most \((n+1)^d\) partitions of \(n\) with at most
\(d\) parts. Hence the series converges absolutely and uniformly.
Its value at coincident eigenvalues follows by continuity.

For the heat coefficients the determinant is an allowed Toeplitz
minor: its row indices are \(i-\lambda_i\), its column indices
are \(j\), and a common integer shift makes them nonnegative.
It is consequently positive. Evaluation at \(U=I\) gives the
displayed sum of probability weights, and a positive mixture
of normalized irreducible characters is of positive type.

We also need the converse at a fixed rank. For an entire real
\(\Psi\) with \(a_0>0\) and \(\Psi(1)=1\), positive type makes
every irreducible coefficient in the preceding expansion
nonnegative, by compact-group orthogonality, equivalently by
\cite[Section~4.1, p.~56]{Faraut2007}. Padding a partition
of length at most \(k\le d\) by zeros multiplies its determinant
by \(a_0^{d-k}\). Thus the specialization \(h_n\mapsto a_n/a_0\)
is nonnegative on every ordinary Schur function of length at most
\(d\).

An allowed minor of order \(k\), with row indices \(i_\ell\)
and column indices \(j_\ell\), becomes the skew Schur value
\(a_0^k s_{\lambda/\mu}\) under this specialization, where
\[
 \lambda_p=j_{k+1-p}+p-i_1-k,\qquad
 \mu_p=i_{k+1-p}+p-i_1-k,\qquad 1\le p\le k.
\]
Transposition and reversal of both lists prove the identity by
ordinary Jacobi--Trudi. The Littlewood--Richardson expansion
of \(s_{\lambda/\mu}\) has nonnegative coefficients and only
partitions contained in \(\lambda\), hence of length at most \(k\).
All these values are nonnegative. Forbidden minors vanish by
upper triangularity. This proves the exact equivalence
\[
 \chi_d\text{ is of positive type on }U(d)
 \quad\Longleftrightarrow\quad (a_k)_{k\ge0}\in PF_d
 \qquad(a_0>0).
\]
It also proves the minimum-order assertion, since \(E_a(-1)>0\)
near zero.

To finish, suppose that all ranks were of positive type for
some \(a>0\). Rank one gives nonnegative Taylor coefficients.
If \(k\) is the first nonzero index, write \(\Psi_a(z)=z^k B(z)\).
The products of \(B\) remain of positive type at every rank,
because they are multiplied by the one-dimensional unitary
character \(\det(U)^{-k}\). Now \(B(0)>0\), so the proved
equivalence makes its coefficient sequence \(PF_\infty\).
Edrei's theorem makes every zero of \(B\), and therefore every
zero of \(\Psi_a\), nonpositive real. But \(H_{-a}\) has a
nonreal zero \(z\) with \(\Re z\ne0\), and
\(1-p(a)z^2\) is a nonreal zero of \(\Psi_a\). This contradiction
proves finiteness of \(d_*(-a)\). Positivity at rank \(d+1\)
restricts to positivity at rank \(d\), since
\(\chi_{a,d+1}(U\oplus1)=\chi_{a,d}(U)\). The intervals already
proved for every fixed \(d\) then give its divergence.
\end{proof}

The formulas specify the full relation with the earlier Toeplitz
test: reciprocal endpoint parameters first undergo
\(z^{-2}\mapsto z^{-2}/p(a)\), and then
\(\beta\mapsto\beta/(1-\beta)\). The character determinant uses
the coefficients after both operations. All claims about the
untranslated detecting order use \(F_a\), as in
Theorem~\ref{thm:fd-limits}.
